\documentclass[11pt]{article}

\usepackage[margin=1in]{geometry}
\usepackage{amsmath,amssymb,amsthm,mathtools}
\usepackage{microtype}
\usepackage{enumitem}
\usepackage{xcolor}
\usepackage{xr-hyper}
\usepackage{hyperref}
\usepackage[nameinlink,noabbrev]{cleveref}

\externaldocument[stack-]{odd_characterization_transport_stack_note}
\externaldocument[delay-]{delayed_odd_control_note}

\hypersetup{
  colorlinks=true,
  linkcolor=blue!60!black,
  citecolor=blue!60!black,
  urlcolor=blue!60!black
}

\theoremstyle{definition}
\newtheorem{definition}{Definition}[section]
\theoremstyle{plain}
\newtheorem{theorem}[definition]{Theorem}
\newtheorem{proposition}[definition]{Proposition}
\newtheorem{corollary}[definition]{Corollary}
\theoremstyle{remark}
\newtheorem{remark}[definition]{Remark}

\title{Selected-Line No-Extra-Collapse Note}
\author{}
\date{}

\begin{document}
\maketitle

Fix an accepted tick \(k\).
This note defines a same-tick selected carrier together with its quotient realization map and
reduced welded image map, and proves that a selected representative with nonzero dominant response
survives both the ambiguity quotient and the reduced welded image on that carrier.

\section{Same-Tick Selected Carrier Data}

\begin{definition}[Same-tick selected-line datum and quotient realization]
\label{def:uom-selected-line-data}
Fix an accepted tick \(k\), and let
\[
\mathfrak D^-_{\Lambda_k}
:=
E^-_{\Lambda_k}/\mathcal A^-_{\Lambda_k},
\qquad
q_{\Lambda_k}:E^-_{\Lambda_k}\to\mathfrak D^-_{\Lambda_k}
\]
be the dominant-response quotient and quotient map of
\cref{stack-def:dominant-response-quotient}.
A \emph{same-tick selected-line datum} at \(\Lambda_k\) consists of:
\begin{enumerate}[leftmargin=2em,label=(\alph*)]
\item a one-dimensional real selected carrier
  \[
  \widehat L^{\mathrm{sel}}_{\Lambda_k}
  =
  \mathbb R\,\widehat\ell_{\Lambda_k};
  \]
\item a selected representative
  \[
  J^{\mathrm{sel}}_{\Lambda_k}\in[\Xi^{(1)}_{\Lambda_k}]\cap E^-_{\Lambda_k}.
  \]
\end{enumerate}
The associated \emph{selected quotient realization map} is
\[
\widehat{\mathfrak R}^{\mathrm{sel}}_{\Lambda_k}
:
\widehat L^{\mathrm{sel}}_{\Lambda_k}\longrightarrow \mathfrak D^-_{\Lambda_k},
\qquad
\widehat{\mathfrak R}^{\mathrm{sel}}_{\Lambda_k}(a\,\widehat\ell_{\Lambda_k})
:=
a\,q_{\Lambda_k}(J^{\mathrm{sel}}_{\Lambda_k}).
\]
\end{definition}

\begin{definition}[Selected reduced welded image and theorem-trivial selected classes]
\label{def:uom-selected-line-trivial}
Fix an accepted tick \(k\) and a same-tick selected-line datum
\cref{def:uom-selected-line-data}.
Let
\[
P_{-,\Lambda_k}^{\mathrm{LoG}}
\]
be the canonical compressed odd projector of
\cref{stack-cor:compressed-export-selection}.
Define the \emph{selected reduced welded image map} by
\[
\mathsf W^{\mathrm{sel}}_{\Lambda_k}
:
\widehat L^{\mathrm{sel}}_{\Lambda_k}\longrightarrow
\mathbb R\,P_{-,\Lambda_k}^{\mathrm{LoG}},
\qquad
\mathsf W^{\mathrm{sel}}_{\Lambda_k}(a\,\widehat\ell_{\Lambda_k})
:=
a\,\rho^{\mathrm{dom}}_{\Lambda_k}(J^{\mathrm{sel}}_{\Lambda_k})\,P_{-,\Lambda_k}^{\mathrm{LoG}}.
\]
The \emph{theorem-trivial selected classes} are
\[
\mathcal T^{\mathrm{inv,sel}}_{\Lambda_k}
:=
\ker\mathsf W^{\mathrm{sel}}_{\Lambda_k}.
\]
\end{definition}

\section{Same-Tick No-Extra-Collapse}

\begin{theorem}[Same-tick no-extra-collapse on the selected line]
\label{thm:uom-selected-line-no-extra-collapse}
Let \(I\) be a connected accepted-branch segment and assume the weak transported descendant-class
interface of \cref{stack-thm:char-weak-interface} on \(I\).
Fix an accepted tick \(k\) with \(\Lambda_k\in I\), and assume a same-tick selected-line datum of
\cref{def:uom-selected-line-data}.
If
\[
\rho^{\mathrm{dom}}_{\Lambda_k}(J^{\mathrm{sel}}_{\Lambda_k})\neq 0,
\]
then
\[
\ker\widehat{\mathfrak R}^{\mathrm{sel}}_{\Lambda_k}
=
\ker\mathsf W^{\mathrm{sel}}_{\Lambda_k}
=
\{0\}.
\]
Equivalently,
\[
\ker\widehat{\mathfrak R}^{\mathrm{sel}}_{\Lambda_k}
=
\mathcal T^{\mathrm{inv,sel}}_{\Lambda_k}
=
\{0\}.
\]
\end{theorem}

\begin{proof}
By \cref{stack-thm:char-weak-interface}, the normalized orientation functional satisfies
\[
\chi_{\Lambda_k}
:=
\nu_{\Lambda_k}^{-1}\rho^{\mathrm{dom}}_{\Lambda_k},
\qquad
\mathcal A^-_{\Lambda_k}\subset\ker\chi_{\Lambda_k}.
\]
Since \(\nu_{\Lambda_k}>0\), this is equivalent to
\[
\mathcal A^-_{\Lambda_k}\subset\ker\rho^{\mathrm{dom}}_{\Lambda_k}.
\]

We first show that the selected quotient point is nonzero.
Assume for contradiction that
\[
q_{\Lambda_k}(J^{\mathrm{sel}}_{\Lambda_k})=0.
\]
Then \(J^{\mathrm{sel}}_{\Lambda_k}\in\mathcal A^-_{\Lambda_k}\), so the kernel inclusion above gives
\[
\rho^{\mathrm{dom}}_{\Lambda_k}(J^{\mathrm{sel}}_{\Lambda_k})=0,
\]
contradicting the nonzero dominant-response hypothesis.
Hence
\[
q_{\Lambda_k}(J^{\mathrm{sel}}_{\Lambda_k})\neq 0.
\]

Now the selected carrier \(\widehat L^{\mathrm{sel}}_{\Lambda_k}\) is one-dimensional and generated by
\(\widehat\ell_{\Lambda_k}\).
The displayed nonvanishing of
\(
q_{\Lambda_k}(J^{\mathrm{sel}}_{\Lambda_k})
\)
implies
\[
\widehat{\mathfrak R}^{\mathrm{sel}}_{\Lambda_k}(\widehat\ell_{\Lambda_k})
=
q_{\Lambda_k}(J^{\mathrm{sel}}_{\Lambda_k})
\neq 0,
\]
and therefore
\[
\ker\widehat{\mathfrak R}^{\mathrm{sel}}_{\Lambda_k}=\{0\}.
\]

By \cref{stack-cor:compressed-export-selection}, the projector
\(
P_{-,\Lambda_k}^{\mathrm{LoG}}
\)
is the canonical rank-one projector onto the compressed odd line and is therefore nonzero.
Since
\(
\rho^{\mathrm{dom}}_{\Lambda_k}(J^{\mathrm{sel}}_{\Lambda_k})\neq 0
\),
we have
\[
\mathsf W^{\mathrm{sel}}_{\Lambda_k}(\widehat\ell_{\Lambda_k})
=
\rho^{\mathrm{dom}}_{\Lambda_k}(J^{\mathrm{sel}}_{\Lambda_k})\,P_{-,\Lambda_k}^{\mathrm{LoG}}
\neq 0,
\]
so
\[
\ker\mathsf W^{\mathrm{sel}}_{\Lambda_k}=\{0\}.
\]
The two kernels are thus equal, and both are trivial.
\end{proof}

\begin{corollary}[Selected representative survives the ambiguity quotient]
\label{cor:uom-selected-line-quotient-nonzero}
In the setting of \cref{thm:uom-selected-line-no-extra-collapse},
\[
q_{\Lambda_k}(J^{\mathrm{sel}}_{\Lambda_k})\neq 0.
\]
\end{corollary}

\begin{proof}
This was proved in the first half of the proof of
\cref{thm:uom-selected-line-no-extra-collapse}.
\end{proof}

\begin{remark}[Scope]
\label{rem:selected-line-no-extra-collapse-scope}
The theorem does not construct the same-tick selected carrier or the selected representative.
It says only that once such data are supplied and the selected representative has nonzero dominant
response, the selected direction survives both the ambiguity quotient and the reduced welded image.
\end{remark}

\begin{remark}[Relation to the delayed selected-line lane]
\label{rem:selected-line-no-extra-collapse-delay}
The same algebraic pattern reappears on the delayed carrier
\cref{delay-def:delayed-selected-line} with delayed realization map
\cref{delay-def:delayed-response-map}; selected-line delayed response sufficiency is proved in
\cref{delay-thm:delayed-response-sufficiency}.
The present note works on the same-tick selected carrier and the current ambiguity quotient,
whereas the delayed note works on the delayed quotient and the next-receiver reduced welded line.
\end{remark}

\end{document}
